titulo: Teorema del valor medio
tema: Derivadas
dificultad: introductoria
etiquetas: derivada, Rolle
actualizada: 2026-07-28
---
\begin{teorema}
Si $f$ es continua en $[a,b]$ y derivable en $(a,b)$, existe $c \in (a,b)$ con
\[
  f'(c) = \frac{f(b)-f(a)}{b-a}.
\]
\end{teorema}

\begin{demostracion}
Definimos
\[
  g(x) = f(x) - \frac{f(b)-f(a)}{b-a}\,(x-a).
\]
Entonces $g$ es continua en $[a,b]$, derivable en $(a,b)$ y $g(a)=g(b)=f(a)$.
Por el teorema de Rolle existe $c \in (a,b)$ tal que $g'(c)=0$, es decir
$f'(c) = \dfrac{f(b)-f(a)}{b-a}$. \qed
\end{demostracion}
